\documentclass[11pt,a4paper]{article}

% ============ PAQUETES ============
\usepackage[utf8]{inputenc}
\usepackage[T1]{fontenc}
\usepackage[spanish]{babel}
\usepackage[margin=2.5cm]{geometry}
\usepackage{graphicx}
\usepackage{tikz}
\usetikzlibrary{shapes,arrows.meta,positioning,automata,calc}
\usepackage{booktabs}
\usepackage{longtable}
\usepackage{array}
\usepackage{xcolor}
\usepackage{enumitem}
\usepackage{fancyhdr}
\usepackage{hyperref}
\usepackage{titlesec}
\usepackage{amssymb}

\hypersetup{
    colorlinks=true,
    linkcolor=blue,
    urlcolor=blue,
    citecolor=black,
    pdftitle={Prueba Practica Unidad IV - ISR-401},
    pdfauthor={Nombre del Estudiante}
}

\setlength{\headheight}{14pt}
\pagestyle{fancy}
\fancyhf{}
\fancyhead[L]{Prueba Práctica -- Unidad IV -- ISR-401}
\fancyhead[R]{\thepage}
\renewcommand{\headrulewidth}{0.4pt}

\titleformat{\section}{\normalfont\Large\bfseries}{\thesection.}{0.6em}{}
\titleformat{\subsection}{\normalfont\large\bfseries}{\thesubsection}{0.6em}{}


\newcommand{\figuraPendiente}[3]{%
  \begin{center}
    \begin{tikzpicture}
      \draw[dashed, thick, gray, fill=gray!8] (0,0) rectangle (#2,#3);
      \node[align=center, text width=#2 cm - 1cm, gray, font=\bfseries] at (#2/2,#3/2) {#1};
    \end{tikzpicture}
  \end{center}%
}

\begin{document}

% =====================================================================
% CARÁTULA
% =====================================================================
\begin{titlepage}
\centering
\vspace*{1cm}

{\large UNIVERSIDAD TÉCNICA ESTATAL DE QUEVEDO\par}
{\normalsize Facultad de Ciencias de la Ingeniería \par}
{\normalsize Carrera de Ingeniería de Software \par}

\vspace{1.2cm}
\rule{0.9\textwidth}{0.4pt}\\[0.4cm]
{\LARGE \bfseries PRUEBA PRÁCTICA -- UNIDAD IV \par}
\vspace{0.3cm}
{\Large Validación, Gestión, Herramientas y Estándares de Requisitos \par}
\rule{0.9\textwidth}{0.4pt}\\[1cm]

{\normalsize Asignatura: Ingeniería de Requisitos (ISR-401) \par}
{\normalsize Docente: Ing. Gleiston Guerrero, Mg. \par}
\vspace{0.6cm}


\begin{tabular}{ll}
\textbf{Estudiante:} & Rizzo Velez Edson Nagib \\[0.4cm]
\textbf{Fecha:} & Martes 11 de Agosto del 2026 \\[0.4cm]
\textbf{Modalidad:} & Individual, en clase \\[0.4cm]
\textbf{Caso:} & Sistema de Gestión de Pedidos \\[0.4cm]
\textbf{Ponderación:} & 100 puntos (30 + 70) \\
\end{tabular}

\vspace{0.8cm}

{\large \textbf{URL del repositorio:} \url{https://github.com/erizzov-boop/Prueba_practica_Unidad_IV_Rizzo}}

\vspace{0.8cm}


\vspace{0.5cm}

\vfill
{\small Documento generado en \LaTeX\ -- compilado a PDF -- versión del repositorio}
\end{titlepage}

% =====================================================================
% CASO PRÁCTICO (contexto, referencia)
% =====================================================================
\section*{Caso práctico: Sistema de Gestión de Pedidos}
Una tienda en línea requiere un módulo para gestionar los pedidos de sus clientes. Un
\textbf{Cliente} (identificado por cédula/RUC, nombre y correo) puede realizar cero o más
\textbf{Pedidos}. Cada \textbf{Pedido} (identificador, fecha y total) contiene una o más
\textbf{Líneas de pedido}, y cada línea referencia un \textbf{Producto} (código, descripción,
precio, stock). Al registrar un pedido, el sistema verifica la disponibilidad de stock: si hay
existencias, confirma el pedido; si no, notifica el faltante. Un pedido confirmado puede
despacharse, luego entregarse, o bien anularse mientras no haya sido entregado. El sistema debe
registrar los pedidos con rapidez y proteger los datos personales del cliente.

\bigskip
\noindent Las diez actividades siguientes usan el mismo dominio; los modelos deben ser
mutuamente consistentes entre sí.

% =====================================================================
% P1
% =====================================================================
\section{Modelo de datos -- Diagrama de clases UML}


\subsection*{Diagrama}
% Diagrama de clases UML - Sistema de Gestión de Pedidos
\begin{tikzpicture}[
    class/.style={
        rectangle split, rectangle split parts=3, draw=black, fill=white,
        text=black, rectangle split part align=left, font=\small,
        text width=4.4cm, rounded corners=1pt
    },
    every edge/.style={draw=black, thick},
    labelfont/.style={font=\scriptsize, text=black}
]

\node[class] (cliente) {
    \textbf{Cliente}
    \nodepart{second} cédulaRUC: String [id] \\ nombre: String \\ correo: String
    \nodepart{third} +registrarPedido()
};

\node[class, right=3.8cm of cliente] (pedido) {
    \textbf{Pedido}
    \nodepart{second} idPedido: String [id] \\ fecha: Date \\ total: Decimal
    \nodepart{third} +confirmar() \\ +despachar() \\ +entregar() \\ +anular()
};

\node[class, below=2.6cm of pedido] (lineapedido) {
    \textbf{LíneaPedido}
    \nodepart{second} cantidad: int \\ subtotal: Decimal
    \nodepart{third} +calcularSubtotal()
};

\node[class, below=2.6cm of cliente] (producto) {
    \textbf{Producto}
    \nodepart{second} código: String [id] \\ descripción: String \\ precio: Decimal \\ stock: int
    \nodepart{third} +verificarStock()
};

% Asociación Cliente - Pedido
\draw (cliente) -- (pedido)
    node[labelfont, pos=0.12, above] {1}
    node[labelfont, pos=0.88, above] {0};

% Composición Pedido - LineaPedido
\draw[-{Diamond[length=8pt,width=6pt]}] (lineapedido) -- (pedido)
    node[labelfont, pos=0.15, right] {1}
    node[labelfont, pos=0.85, right] {1};

% Asociación LineaPedido - Producto
\draw (producto) -- (lineapedido)
    node[labelfont, pos=0.12, above] {1}
    node[labelfont, pos=0.88, above] {0};

\end{tikzpicture}








% =====================================================================
% P2
% =====================================================================
\section{Modelo funcional -- Diagrama de actividades UML}

\subsection*{Diagrama}
% Diagrama de actividades - Registrar pedido
\begin{tikzpicture}[
    node distance=1.1cm and 2.2cm,
    initial/.style={circle, draw=black, fill=black, minimum size=0.35cm},
    final/.style={circle, draw=black, fill=white, minimum size=0.45cm, double, double distance=1.3pt},
    action/.style={rectangle, rounded corners=5pt, draw=black, fill=white, text=black,
        minimum width=3.2cm, minimum height=0.9cm, align=center, font=\small},
    decision/.style={diamond, draw=black, fill=white, text=black, aspect=2,
        inner sep=1pt, align=center, font=\scriptsize},
    every path/.style={draw=black, thick, -{Latex[length=2.5mm]}},
    lbl/.style={font=\scriptsize, text=black}
]

\node[initial] (start) {};
\node[action, below=of start] (recibir) {Recibir pedido};
\node[decision, below=1.3cm of recibir] (stock) {¿Stock\\disponible?};
\node[action, below left=1.4cm and 0.3cm of stock] (confirmar) {Confirmar pedido};
\node[action, below right=1.4cm and 0.3cm of stock] (notificar) {Notificar faltante};
\node[final, below=2.6cm of stock] (end) {};

\draw (start) -- (recibir);
\draw (recibir) -- (stock);
\draw (stock) -| (confirmar) node[lbl, pos=0.25, above] {Sí};
\draw (stock) -| (notificar) node[lbl, pos=0.25, above] {No};
\draw (confirmar) |- (end);
\draw (notificar) |- (end);

\end{tikzpicture}









% =====================================================================
% P3
% =====================================================================
\section{Modelo de comportamiento -- Máquina de estados UML}
\subsection*{Diagrama}
\begin{figure}[h!]
    \centering
    \includegraphics[width=0.85\linewidth]{DiagramaEstados.png}
    \caption{Diagrama de estados}
    \label{fig:placeholder}
\end{figure}

% =====================================================================
% P4
% =====================================================================
\section{Consistencia entre las tres perspectivas}

\subsection*{Tabla de verificación cruzada}
\begin{table}[h!]
\centering
\begin{tabular}{lll}
\toprule
\textbf{Evento (P3)} & \textbf{Función/acción (P2)} & \textbf{Entidad/atributo (P1)} \\
\midrule
confirmar  & Confirmar pedido & Pedido.estado \\
despachar  & Despachar pedido & Pedido.estado \\
entregar   & Entregar pedido  & Pedido.estado \\
anular     & Anular pedido    & Pedido.estado \\
\bottomrule
\end{tabular}
\caption{Verificación cruzada entre eventos, funciones y entidades}
\end{table}

\subsection*{Inconsistencia detectada}
El evento \textbf{``notificar faltante''} (P2) no tiene una transición correspondiente en la
máquina de estados de P3, ya que el pedido permanece en el estado \emph{Registrado}.

\subsection*{Corrección aplicada}
Se documenta como nota en la máquina de estados que el camino no produce cambio de estado.

% =====================================================================
% P5
% =====================================================================
\section{Especificación de requisitos con esquema de atributos}

\begin{description}[leftmargin=1.5cm]
  \item[RF-01:] El sistema debe permitir registrar un pedido con sus líneas y calcular el total
        automáticamente.
  \item[RF-02:] El sistema debe verificar el stock del producto antes de confirmar el pedido.
  \item[RNF-01:] \textit{(Eficiencia de desempeño, ISO/IEC 25010)}: el sistema debe registrar un
        pedido en menos de 3 segundos bajo carga normal.
  \item[RNF-02:] \textit{(Seguridad, ISO/IEC 25010)}: los datos personales del cliente deben
        cifrarse en reposo y en tránsito.
\end{description}

\subsection*{Esquema de atributos}
\begin{table}[h!]
\centering
\begin{tabular}{lllllll}
\toprule
\textbf{ID} & \textbf{Prioridad} & \textbf{Estado} & \textbf{Fuente} & \textbf{Estabilidad} & \textbf{Versión} & \textbf{Verificación} \\
\midrule
RF-01  & Alta  & Propuesto & Caso de negocio & Alta  & 1.0 & Prueba \\
RF-02  & Alta  & Propuesto & Caso de negocio & Alta  & 1.0 & Prueba \\
RNF-01 & Media & Propuesto & Caso de negocio & Media & 1.0 & Prueba de carga \\
RNF-02 & Media & Propuesto & Caso de negocio & Media & 1.0 & Inspección/prueba \\
\bottomrule
\end{tabular}
\caption{Esquema de atributos de los requisitos}
\end{table}

% =====================================================================
% P6
% =====================================================================
\section{Priorización MoSCoW}

\begin{itemize}[leftmargin=1.2cm]
  \item \textbf{RF-01 -- Must have}: Es el núcleo funcional del sistema.
  \item \textbf{RF-02 -- Must have}: Garantiza la integridad de los datos e inventario.
  \item \textbf{RNF-01 -- Should have}: Mejora la experiencia pero no bloquea la operación básica.
  \item \textbf{RNF-02 -- Could have}: Deseable desde el inicio puede incorporarse en una iteración posterior.
\end{itemize}

% =====================================================================
% P7
% =====================================================================
\section{Validación por inspección (checklist ISO/IEC/IEEE 29148)}

\subsection*{Propiedades evaluadas}
No ambiguo, completo, singular, verificable, factible.

\subsection*{Defectos detectados (sin corregir)}
RNF-02 no especifica algoritmo ni umbral medible
$\Rightarrow$ incumple la propiedad de \emph{``verificable''}.

\subsection*{Retrabajo (segundo paso)}
``Los datos personales deben cifrarse con AES-256 en reposo y TLS 1.3 en tránsito.''

% =====================================================================
% P8
% =====================================================================
\section{Pruebas de aceptación trazadas}

\begin{table}[h!]
\centering
\small
\begin{tabular}{lllll}
\toprule
\textbf{ID} & \textbf{Tipo} & \textbf{Entradas} & \textbf{Criterio de éxito} & \textbf{Requisito que traza} \\
\midrule
PA-01 & Funcional     & Pedido con 2 líneas válidas    & Total calculado correctamente          & RF-01 \\
PA-02 & Funcional     & Producto sin stock             & Notifica faltante, no confirma pedido  & RF-02 \\
PA-03 & No funcional  & 100 pedidos simultáneos        & Cada registro $<$ 3 sg                  & RNF-01 \\
PA-04 & No funcional  & Intento de acceso directo a BD & Datos ilegibles sin clave              & RNF-02 \\
\bottomrule
\end{tabular}
\caption{Pruebas de aceptación trazadas a los requisitos}
\end{table}

% =====================================================================
% P9
% =====================================================================
\section{Matriz de trazabilidad}

\begin{table}[h!]
\centering
\small
\begin{tabular}{lp{2.3cm}p{5.2cm}l}
\toprule
\textbf{Requisito} & \textbf{Fuente} & \textbf{Modelo asociado} & \textbf{Prueba de aceptacion} \\
\midrule
RF-01  & Caso de negocio & P1: Pedido/LíneaPedido; P2: flujo completo & PA-01 \\
RF-02  & Caso de negocio & P1: Producto.stock; P2: decisión stock     & PA-02 \\
RNF-01 & Caso de negocio & P2: flujo ``Registrar pedido''             & PA-03 \\
RNF-02 & Caso de negocio & P1: Cliente (datos personales)             & PA-04 \\
\bottomrule
\end{tabular}
\caption{Matriz de trazabilidad de requisitos}
\end{table}

\subsection*{Dependencia horizontal}
RF-02 depende de RF-01, ya que no puede verificarse el stock sin que exista previamente un
pedido creado.

\subsection*{Requisitos huérfanos}
Ninguno.

% =====================================================================
% P10
% =====================================================================
\section{Gestión del cambio y línea base}

\subsection*{Solicitud de cambio}
``Agregar notificación por correo al cliente al confirmar el pedido''.

\subsection*{Análisis de impacto (según matriz de P9)}
Afecta a RF-01, al diagrama de clases (nuevo método en la clase Pedido) y a la prueba PA-01.

\subsection*{Decisión del CCB}
\textbf{Aprobada}, con prioridad \emph{Should have}.

\subsection*{Nueva versión}
RF-01 pasa de la versión 1.0 a la versión 1.1.

\subsection*{Línea base declarada (Baseline 1.0)}
Se congela la siguiente configuración coherente:
\begin{itemize}[leftmargin=1.2cm]
  \item ERS V1.0
  \item Modelos de clases, actividades y estados V1.0
  \item Matriz de trazabilidad V1.0
  \item Pruebas de aceptación V1.0
\end{itemize}


\end{document}
